
Pour cela, on considère le complexe :
@@M0@@
où @@M1@@ est l'algèbre symétrique de @@M2@@. C'est une résolution de @@M3@@ sur @@M4@@ (VI 1.11). Le complexe associé sur @@M5@@, noté @@M6@@, a pour @@M7@@-ième terme @@M8@@ et est donc une résolution de @@M9@@ par des @@M10@@-Modules localement libres de type fini ; on a donc @@M11@@. Par suite :
@@M12@@
La proposition en résulte par additivité.

\subsection*{2. Fibrés principaux sous les tores déployés}

2.1. Soient @@M13@@ un schéma noethérien, et @@M14@@ un tore sur @@M15@@. On suppose @@M16@@ déployé; on a donc un isomorphisme @@M17@@, où @@M18@@ est un groupe abélien libre de rang @@M19@@, soit encore @@M20@@. Soit @@M21@@ un torseur sous @@M22@@. Alors @@M23@@ est affine sur @@M24@@, défini par une @@M25@@-Algèbre @@M26@@ graduée de type @@M27@@ ; si on pose @@M28@@, alors @@M29@@ est un @@M30@@-Module inversible pour tout @@M31@@, et, pour @@M32@@, l'homomorphisme @@M33@@ défini par la multiplication de @@M34@@ est un isomorphisme (SGA 3 VIII 4.1). On définit donc un homomorphisme multiplicatif @@M35@@ en associant à @@M36@@ la classe de @@M37@@ dans @@M38@@.

